\documentclass[11pt,a4paper]{amsart}
\usepackage{amsmath,amssymb,amsthm}
\usepackage[margin=1in]{geometry}
\usepackage[hidelinks]{hyperref}
\setlength{\emergencystretch}{2em}

\newtheorem*{theorem}{Theorem}
\newtheorem{lemma}{Lemma}
\theoremstyle{remark}
\newtheorem*{remark}{Remark}

\title[Erd\H{o}s Problem 878, Question 3b]{Erd\H{o}s Problem 878, Question 3b:\\ the maxima of $f$ and $F$ differ infinitely often}
\author{Ilia Pauzner}
\date{}

\begin{document}
\maketitle

\section{The problem}

Throughout, $p$ denotes a prime. For an integer $n \ge 1$ let $\mathcal{P}(n)$ be the set of primes dividing $n$, let $\omega(n) = |\mathcal{P}(n)|$, and let
\[
q_p(n) = p^{\lfloor \log_p n \rfloor}
\]
be the largest power of $p$ that does not exceed $n$. Erd\H{o}s~\cite{Er84e} considered the functions
\[
f(n) = \sum_{p \in \mathcal{P}(n)} q_p(n)
\qquad\text{and}\qquad
F(n) = \max_A \sum_{a \in A} a,
\]
where the maximum is taken over all sets $A \subseteq \{2, 3, \dots, n\}$ of pairwise coprime integers all of whose prime factors divide $n$; we call such sets \emph{admissible for $n$}. The set $\{q_p(n) : p \in \mathcal{P}(n)\}$ is admissible, so $F(n) \ge f(n)$. For real $x \ge 1$ put
\[
M_f(x) = \max_{1 \le n \le x} f(n), \qquad M_F(x) = \max_{1 \le n \le x} F(n);
\]
thus $M_F(x) \ge M_f(x)$, and both functions are non-decreasing. Erd\H{o}s asked whether $M_f(x) = M_F(x)$ for all sufficiently large $x$~\cite{EP878}. The answer is no.

\begin{theorem}
For every $X_0$ there is an integer $x \ge X_0$ with $M_f(x) < M_F(x)$. In particular, $M_f(x) = M_F(x)$ fails for arbitrarily large $x$.
\end{theorem}

Lower bounds for $F$ will always come from the following observation.

\begin{remark}[Building admissible sets]
Let $a_1, \dots, a_k \in \{2, \dots, n\}$, and suppose that their sets of prime factors are pairwise disjoint subsets of $\mathcal{P}(n)$. Then $a_1, \dots, a_k$ are pairwise coprime, hence distinct, so $\{a_1, \dots, a_k\}$ is admissible for $n$ and $a_1 + \dots + a_k \le F(n)$.
\end{remark}

\section{Idea of the proof}

Suppose that $M_F(x) = M_f(x)$ for all large $x$, and let $n \le y$ be a point where $f$ attains its maximum $M_f(y)$ on $[1, y]$. The value $f(n)$ is a sum of prime powers, each at most $n$, one for each prime factor of $n$. We construct a number $m$ only slightly larger than $n$ with $F(m) \ge f(n) + Y/25$, where $Y \le y \le 1.05\,Y$ for a fixed scale $Y$:
\begin{itemize}
\item $m$ is divisible by almost all prime factors of $n$, so $F(m)$ can reuse the corresponding terms of $f(n)$;
\item two large prime factors $r_a, r_b$ of $n$ are dropped, which loses at most $2n$;
\item six fixed primes $s_1, \dots, s_6$ divide $m$. The scale $Y$ is chosen so that their terms in $f(n)$ are small (less than $\frac{6}{7}n$ in total), and in $F(m)$ we replace them by three numbers $b_1, b_2, b_3 \in (0.98\,m, m]$, each built from two of the $s_i$, which contribute more than $2.94\,m$.
\end{itemize}
Since $2.94 > 2 + \frac{6}{7}$, this gains a fixed proportion of $n$. As $M_F = M_f$, the function $M_f$ must therefore increase by at least $Y/25$ on the short interval from $y$ to $y + (m - n)$. If $r_a, r_b \ge M$, the step $m - n$ is at most $2CY/M^2$, where $C = s_1 \cdots s_6$; and $n$ has more than $M$ prime factors as soon as $M_f(y)/y$ is large compared with $M$. Iterating, $M_f$ grows faster and faster while the steps shrink like $1/M^2$, so infinitely many jumps fit into the bounded window $[Y, 1.05\,Y]$, which is impossible.

\section{Four lemmas}

\begin{lemma}[Smooth numbers just below $n$]\label{lem:smooth}
Let $p < r$ be primes with $98 r < 100 p$. For every $n \ge p^{p^2 + 2}$ there is an integer $z$ with $0.98\, n < z \le n$ all of whose prime factors lie in $\{p, r\}$; moreover $z \ge 2$.
\end{lemma}

\begin{proof}
Let $\ell = \lfloor \log_p n \rfloor$, so that $p^\ell \le n < p^{\ell + 1}$ and $\ell \ge p^2 + 2$. Put $J = p^2 + 1 < \ell$ and
\[
z_j = p^{\ell - j} r^j \qquad (0 \le j \le J).
\]
Each term of this sequence is $r/p < 100/98$ times the previous one. It starts at $z_0 = p^\ell \le n$ and ends above $n$: since $r \ge p + 1$,
\[
z_J \ge p^{\ell - J} (p + 1)^J \ge p^{\ell - J} \bigl(p^J + J p^{J - 1}\bigr) = p^\ell + p^{\ell + 1} + p^{\ell - 1} > p^{\ell + 1} > n.
\]
Let $j < J$ be the last index with $z_j \le n$. Then $n < z_{j + 1} = z_j \cdot r/p < \frac{100}{98} z_j$, so $z = z_j$ satisfies $0.98\, n < z \le n$. Also $z \ge z_0 \ge p \ge 2$.
\end{proof}

\begin{lemma}[Maximizers lie in the upper half]\label{lem:double}
For every $n \ge 1$ we have $f(2n) > f(n)$. Consequently, if $1 \le n \le y$ and $f(n) = M_f(y)$, then $n > y/2$.
\end{lemma}

\begin{proof}
Every prime factor $p$ of $n$ divides $2n$, and $q_p(2n) \ge q_p(n)$. The term for $p = 2$ strictly increases: if $2 \mid n$, then $q_2(2n) = 2 q_2(n)$, and if $2 \nmid n$, then $f(2n)$ has the additional term $q_2(2n) \ge 2$. Hence $f(2n) > f(n)$. If $n \le y/2$, then $2n \le y$ and $f(2n) > f(n) = M_f(y)$, which is impossible.
\end{proof}

\begin{lemma}[$M_f$ grows faster than linearly]\label{lem:superlinear}
For every $\lambda > 0$ we have $M_f(y) > \lambda y$ for all sufficiently large $y$.
\end{lemma}

\begin{proof}
Since $\sum_p 1/p$ diverges, there is $N$ with $\sum_{p < N} 1/p \ge 2\lambda + 1$. Let $P = \prod_{p < N} p$ and $y \ge 2P$. The number $t = P \lfloor y/P \rfloor$ satisfies $y/2 < t \le y$ (because $t > y - P \ge y/2$) and is divisible by every prime $p < N$, and $q_p(t) > t/p$ for each such $p$. Hence
\[
M_f(y) \ge f(t) \ge \sum_{p < N} q_p(t) > t \sum_{p < N} \frac{1}{p} \ge (2\lambda + 1)\, t > (2\lambda + 1) \frac{y}{2} > \lambda y. \qedhere
\]
\end{proof}

\medskip
\noindent\textbf{Six special primes.} Let
\[
\begin{aligned}
s_1 &= 168000061, & s_2 &= 168000101, & s_3 &= 168000121, \\
s_4 &= 168000139, & s_5 &= 168000143, & s_6 &= 168000167.
\end{aligned}
\]
These six numbers are primes with $s_1 < s_2 < \dots < s_6$, and apart from this the proof uses only the following two properties:
\begin{itemize}
\item $98\, s_{2i} < 100\, s_{2i - 1}$ for $i = 1, 2, 3$ (the primes differ by less than $200$);
\item every $s_i$ is larger than $14 \cdot \frac{21}{20} \cdot 15^6 = 167\,442\,187.5$.
\end{itemize}
Let $S = \{s_1, \dots, s_6\}$ and $C = s_1 s_2 \cdots s_6$.

\begin{lemma}[Choosing the scale]\label{lem:scale}
Call an integer $Y \ge 1$ \emph{good} if $q_p(\lfloor \tfrac{21}{20} Y \rfloor) < \tfrac{1}{14} Y$ for every $p \in S$. For every integer $Y_0 \ge 1$, one of the seven numbers $Y_0, 15 Y_0, 15^2 Y_0, \dots, 15^6 Y_0$ is good.
\end{lemma}

\begin{proof}
Fix $p \in S$; we show that the inequality for $p$ fails for at most one of the numbers $Y_i = 15^i Y_0$. Suppose it fails for $Y_i$ and $Y_j$ with $i < j$, and let $u = q_p(\lfloor \frac{21}{20} Y_i \rfloor)$ and $v = q_p(\lfloor \frac{21}{20} Y_j \rfloor)$, so that $u \ge \frac{1}{14} Y_i$ and $v \ge \frac{1}{14} Y_j$. Since $Y_j \ge 15\, Y_i$,
\[
u \le \tfrac{21}{20} Y_i < \tfrac{15}{14} Y_i \le \tfrac{1}{14} Y_j \le v.
\]
Thus $v > u$, and as both are powers of $p$, $v \ge p\, u \ge \frac{p}{14} Y_i$. On the other hand, $v \le \frac{21}{20} Y_j \le \frac{21}{20} \cdot 15^6\, Y_i$. Together these give $p \le 14 \cdot \frac{21}{20} \cdot 15^6$, contrary to the choice of the $s_i$. Since $S$ has six elements and there are seven numbers $Y_i$, some $Y_i$ satisfies the inequality for every $p \in S$.
\end{proof}

\section{The key construction}

Let $N_0 = s_5^{\,s_5^2 + 2}$. Since $s_1 < s_3 < s_5$, Lemma~\ref{lem:smooth} applies to each of the pairs $(s_1, s_2)$, $(s_3, s_4)$, $(s_5, s_6)$ for every $n \ge N_0$.

\begin{lemma}[Three pairs]\label{lem:pairs}
Let $m \ge N_0$ be divisible by $C$, and let $R$ be a set of primes dividing $m$ with $R \cap S = \emptyset$. Then
\[
F(m) > 2.94\, m + \sum_{p \in R} q_p(m).
\]
\end{lemma}

\begin{proof}
By Lemma~\ref{lem:smooth}, for $i = 1, 2, 3$ there is an integer $b_i$ with $0.98\, m < b_i \le m$ whose prime factors lie in $\{s_{2i - 1}, s_{2i}\}$. The numbers $b_1, b_2, b_3$ and $q_p(m)$ for $p \in R$ lie in $\{2, \dots, m\}$, and their sets of prime factors are contained in $\{s_1, s_2\}$, $\{s_3, s_4\}$, $\{s_5, s_6\}$, and $\{p\}$ for $p \in R$, respectively. These are pairwise disjoint subsets of $\mathcal{P}(m)$, so by the Remark
\[
F(m) \ge b_1 + b_2 + b_3 + \sum_{p \in R} q_p(m) > 3 \cdot 0.98\, m + \sum_{p \in R} q_p(m). \qedhere
\]
\end{proof}

\section{Proof of the theorem}

Assume, for a contradiction, that there is $X_0$ with
\begin{equation}\label{eq:assumption}
M_F(x) \le M_f(x) \qquad \text{for all integers } x \ge X_0.
\end{equation}
Put $B = 2000\, C + 1$. Fix a good integer $Y$ such that
\begin{equation}\label{eq:Y}
Y \ge X_0, \qquad Y \ge 2 N_0, \qquad M_f(Y) > 2 (B + 10)\, Y;
\end{equation}
such $Y$ exists by Lemma~\ref{lem:scale} applied to an integer $Y_0$ larger than $X_0$, $2 N_0$, and the threshold given by Lemma~\ref{lem:superlinear} for $\lambda = 2(B + 10)$.

\begin{lemma}[One jump]\label{lem:jump}
Let $M \ge 1$ and $y$ be integers with $Y \le y \le \frac{21}{20} Y$ and $M_f(y) \ge \frac{21}{20} (M + 8)\, Y$. Then there is an integer $y'$ with
\[
y \le y' \le y + \frac{2 C Y}{M^2} \qquad\text{and}\qquad M_f(y') > M_f(y) + \frac{Y}{25}.
\]
\end{lemma}

\begin{proof}
Choose $n \le y$ with $f(n) = M_f(y)$. By Lemma~\ref{lem:double}, $n > y/2 \ge Y/2$; hence $Y < 2n$ and, by~\eqref{eq:Y}, $n > N_0$. Also $n \le \lfloor \frac{21}{20} Y \rfloor \le 2Y$.

\emph{Many prime factors.} Every term of $f(n)$ is at most $n$, so
\[
\tfrac{21}{20} (M + 8)\, Y \le f(n) \le \omega(n)\, n \le \tfrac{21}{20}\, \omega(n)\, Y,
\]
and therefore $\omega(n) \ge M + 8$.

\emph{Two large primes.} At most $6$ prime factors of $n$ lie in $S$, and fewer than $M$ primes are smaller than $M$. Hence $n$ has two distinct prime factors $r_a, r_b \notin S$ with $r_a \ge M$ and $r_b \ge M$.

\emph{The number $m$.} Let $R = \mathcal{P}(n) \setminus (S \cup \{r_a, r_b\})$ be the set of the remaining prime factors of $n$, and let $D = C \prod_{p \in R} p$. Since $r_a r_b \prod_{p \in R} p$ divides $n$,
\[
\prod_{p \in R} p \le \frac{n}{r_a r_b} \le \frac{2Y}{M^2}, \qquad\text{so}\qquad D \le \frac{2 C Y}{M^2}.
\]
Let $m$ be the smallest multiple of $D$ that is larger than $n$. Then $n < m \le n + D$, and every prime in $R \cup S$ divides $m$.

\emph{Upper bound for $f(n)$.} The prime factors of $n$ are the primes in $R$, the primes $r_a$ and $r_b$, and possibly some primes in $S$. For $p \in S$, since $Y$ is good and $n \le \lfloor \frac{21}{20} Y \rfloor$, we have $q_p(n) \le q_p(\lfloor \frac{21}{20} Y \rfloor) < \frac{Y}{14} < \frac{n}{7}$. For $p \in \{r_a, r_b\}$ we use $q_p(n) \le n$, and for $p \in R$ we use $q_p(n) \le q_p(m)$. Hence
\[
f(n) < \frac{6}{7}\, n + 2n + \sum_{p \in R} q_p(m) = \frac{20}{7}\, n + \sum_{p \in R} q_p(m).
\]

\emph{Lower bound for $F(m)$.} As $m > n > N_0$, Lemma~\ref{lem:pairs} applies to $m$ and $R$ and gives
\[
F(m) > 2.94\, m + \sum_{p \in R} q_p(m) > 2.94\, n + \sum_{p \in R} q_p(m).
\]

\emph{Comparison.} Subtracting the two bounds and using $n > Y/2$,
\[
F(m) - f(n) > \Bigl(2.94 - \frac{20}{7}\Bigr) n = \frac{29}{350}\, n > \frac{29}{700}\, Y > \frac{Y}{25}.
\]
Now let $y' = y + D$. Then $y \le y' \le y + \frac{2CY}{M^2}$, $m \le y'$, and $y' \ge Y \ge X_0$, so by~\eqref{eq:assumption}
\[
M_f(y') \ge M_F(y') \ge F(m) > f(n) + \frac{Y}{25} = M_f(y) + \frac{Y}{25}. \qedhere
\]
\end{proof}

\medskip
\noindent\textbf{Iterating the jump.} We construct integers $y_0 \le y_1 \le y_2 \le \cdots$ such that for every $k \ge 0$
\begin{equation}\label{eq:invariant}
Y \le y_k \le Y + 54\, C\, Y \Bigl(\frac{1}{B} - \frac{1}{B + k}\Bigr)
\qquad\text{and}\qquad
M_f(y_k) \ge \frac{21}{20} (B + 9 + k)\, Y.
\end{equation}
Since $B > 1080\, C$, we have $54\, C\, Y / B < Y/20$, so all points considered below lie in the window $[Y, \frac{21}{20} Y]$.

For $k = 0$ take $y_0 = Y$; the second condition holds because $M_f(Y) > 2(B + 10)\, Y \ge \frac{21}{20}(B + 9)\, Y$.

Suppose $y_k$ satisfies~\eqref{eq:invariant}. Put $M = B + 1 + k$, so that $\frac{21}{20}(M + 8)\, Y = \frac{21}{20}(B + 9 + k)\, Y$, and apply Lemma~\ref{lem:jump} $27$ times in succession, starting from $y_k$; let $y_{k+1}$ be the last point obtained. Each application is legitimate. Indeed, $M_f$ is non-decreasing and the points increase, so the condition $M_f(y) \ge \frac{21}{20}(M + 8)\, Y$ persists. Moreover, after $j \le 27$ jumps the point is at most $y_k + j \cdot 2CY/M^2$, so the total distance travelled is at most
\[
27 \cdot \frac{2 C Y}{(B + 1 + k)^2} \le \frac{54\, C\, Y}{(B + k)(B + k + 1)} = 54\, C\, Y \Bigl(\frac{1}{B + k} - \frac{1}{B + k + 1}\Bigr),
\]
so every point stays at most $Y + 54\, C\, Y \bigl(\frac{1}{B} - \frac{1}{B + k + 1}\bigr) < \frac{21}{20} Y$. In particular, $y_{k+1}$ satisfies the first condition of~\eqref{eq:invariant}. Since each jump increases $M_f$ by more than $Y/25$,
\[
M_f(y_{k + 1}) > M_f(y_k) + \frac{27}{25}\, Y \ge \frac{21}{20} (B + 9 + k)\, Y + \frac{21}{20}\, Y = \frac{21}{20} (B + 10 + k)\, Y,
\]
which is the second condition for $k + 1$.

\medskip
\noindent\textbf{The contradiction.} By~\eqref{eq:invariant}, every $y_k$ is at most $\frac{21}{20} Y$, so $M_f(y_k) \le M_f(\frac{21}{20} Y)$, a number that does not depend on $k$. On the other hand, $M_f(y_k) \ge \frac{21}{20}(B + 9 + k)\, Y \ge k$ for every $k$. Taking $k > M_f(\frac{21}{20} Y)$ gives a contradiction, which proves the Theorem. \qed

\begin{remark}[How the constants fit together]
The numbers in the proof are tied together as follows.
\begin{itemize}
\item The gain in one jump comes from $3 \cdot 0.98 = 2.94 > 2 + \frac{6}{7}$, and $\bigl(2.94 - \frac{20}{7}\bigr) \cdot \frac{1}{2} = \frac{29}{700} > \frac{1}{25}$.
\item $27$ jumps raise $M_f$ by more than $\frac{27}{25} Y > \frac{21}{20} Y$, which is enough to raise $M$ by one.
\item The steps have length at most $2CY/M^2$ and $\sum_M 1/M^2$ converges. Starting from $M = B + 1$ with $B > 1080\, C$, the total distance is at most $54\, C\, Y / B < Y/20$, so we never leave the window $[Y, \frac{21}{20} Y]$, on which the goodness of $Y$ controls the terms of the six special primes.
\item The bound $s_i > 14 \cdot \frac{21}{20} \cdot 15^6$ makes the pigeonhole argument of Lemma~\ref{lem:scale} work, and $98\, s_{2i} < 100\, s_{2i-1}$ is what Lemma~\ref{lem:smooth} needs.
\end{itemize}
The same argument is formalized in Lean in \nolinkurl{FormalConjectures/ErdosProblems/878.lean} (theorem \texttt{erdos\_878\_q3\_disproved}).
\end{remark}

\begin{thebibliography}{9}
\bibitem{Er84e}
P.~Erd\H{o}s, \emph{On two unconventional number theoretic functions and on some related problems}, 1984.

\bibitem{EP878}
\emph{Erd\H{o}s Problem \#878}, \url{https://www.erdosproblems.com/878}.
\end{thebibliography}

\end{document}
